
\section{\BookRecordIndex{T0120}非挤压定理}\label{sec:9-3}

\personlater{geluomofu}{格罗莫夫}{Gromov}的\BookRecordIndex{T0120}非挤压定理是辛拓扑的一块基石。它断言，
辛同胚不可能把一个球挤进半径更小的柱体。用 $B^{2n}(r)$ 表示
$\mathbb R^{2n}$ 中以原点为中心、半径为 $r$ 的闭球。

\begin{theorem}[格罗莫夫]\label{thm:9-3-1}
若 $\iota:B^{2n}(r)\to\mathbb R^{2n}$ 是满足
\[
  \iota\bigl(B^{2n}(r)\bigr)\subset B^2(R)\times\mathbb R^{2n-2}
\]
的辛嵌入，则 $r\le R$。
\end{theorem}

反过来，这一性质刻画了辛同胚与反辛同胚；该结论由 Eliashberg \srcite{98} 和
Ekeland--Hofer \srcite{97} 独立证明。由此得到辛刚性的一种体现：辛同胚群关于
$C^0$ 拓扑是闭的。这个现象在 \srcite{277} 第 12 章中已有完整讨论，故此处不再展开。
我们转而集中给出\BookRecordIndex{T0120}非挤压定理如下推广的完整证明。

\begin{theorem}\label{thm:9-3-2}
设 $(V,\tau)$ 是维数为 $2n-2$ 的紧\BookRecordIndex{T0358}辛流形，且 $\pi_2(V)=0$。若球
$(B^{2n}(r),\omega_0)$ 可以辛嵌入 $B^2(R)\times V$，则 $r\le R$。
\end{theorem}

定理~\ref{thm:9-3-2} 蕴含定理~\ref{thm:9-3-1}。事实上，一个辛嵌入
\[
  B^{2n}(r)\lhook\joinrel\longrightarrow B^2(R)\times\mathbb R^{2n-2}
\]
会诱导辛嵌入
$B^{2n}(r)\lhook\joinrel\longrightarrow B^2(R)\times V$，其中
\[
  V:=\mathbb R^{2n-2}/\lambda\mathbb Z^{2n-2}
\]
是 $(2n-2)$-维环面，而 $\lambda>0$ 取得充分大。因此由定理~\ref{thm:9-3-2}，
$R\ge r$，这就证明了定理~\ref{thm:9-3-1}。

为证明定理~\ref{thm:9-3-2}，把二维圆盘 $B^2(R)$ 换成二维球面。论证建立在如下观察
之上：对乘积 $M:=S^2\times V$ 和类 $A:=[S^2\times\mathrm{pt}]$，\BookRecordIndex{T0143}格罗莫夫--威滕
不变量为
\[
  \operatorname{GW}^{M}_{A,1}\bigl(\operatorname{PD}(\mathrm{pt})\bigr)=1.
\]
由于\BookRecordIndex{T0246}模空间 $\mathcal M(A;J)/G$ 紧，这意味着：对 $M$ 上每个被乘积\BookRecordIndex{T0363}辛形式\BookRecordIndex{T0372}驯服的\BookRecordIndex{T0088}殆复
结构，以及每个点 $x_0\in M$，都必有一条通过 $x_0$ 的 $J$-全纯 $A$-球面。特别地，
当 $\iota^*J=J_0$ 且 $x_0=\iota(0)$ 时亦然。在这个特殊情形，相应的 \BookRecordIndex{T0003}$J$-全纯曲线
拉回到 $B^{2n}(r)$ 中一张通过原点的\BookRecordIndex{T0175}极小曲面。\BookRecordIndex{T0092}单调性定理断言，这张曲面的面积至少为
$\pi r^2$。因而，类 $A$ 的辛面积（它等于该二维球面的面积）必至少为 $\pi r^2$。
下面将改造这个论证，以辛吹起构造绕开\BookRecordIndex{T0175}极小曲面\BookRecordIndex{T0092}单调性定理。Lalonde--Pinsonnault
\srcite{229} 中也有类似论证。该文还对\BookRecordIndex{T0137}高阶宽度作了有趣的讨论：给定
$M$ 的辛同胚群第 $k$ 个同伦群中的非平凡元 $\alpha$，他们考察这样一种辛球的最大半径：
在 $\alpha$ 的某个代表元中，所有辛同胚都固定该球的像。

\paragraph{吹起构造。}

在点 $x_0$ 处吹起复流形 $(M,J)$ 的朴素想法，是用通过 $x_0$ 的所有复直线之集替代
$x_0$。为了把这个定义说得更正式，考虑 $\mathbb{CP}^{n-1}$ 上\BookRecordIndex{T0414}重言线丛的全空间
\[
  L:=\bigl\{(z,[w_1:\cdots:w_n])\in\mathbb C^n\times\mathbb{CP}^{n-1}
  \mid z\in\mathbb Cw\bigr\}.
\]
映射 $(z,[w_1:\cdots:w_n])\mapsto z$ 把 $L$ 中零截面的补集与
$\mathbb C^n\setminus\{0\}$ 识别。因此，$L$ 可以看成 $\mathbb C^n$ 的一个副本，其中
原点 $0$ 已被 $L$ 的零截面替代。零截面正是 $\mathbb C^n$ 中通过 $0$ 的所有直线之集，
故可把 $L$ 看作 $\mathbb C^n$ 在 $0$ 处吹起的模型。

一般复流形的吹起如下进行。选取复嵌入
$\iota:B^{2n}(\delta)\to M$，使 $\iota(0)=x_0$，并对
$0<r<\delta$ 考虑 $L$ 的 $r$-球子丛\sourcefootnote{原书定义嵌入时只给出半径
$\delta$，随后直接使用 $U_r$，却漏写了保证 $\iota(z)$ 有定义所必需的
$0<r<\delta$。}
\[
  U_r:=\bigl\{(z,[w_1:\cdots:w_n])\in\mathbb C^n\times\mathbb{CP}^{n-1}
  \mid z\in\mathbb Cw,\ |z|<r\bigr\}.
\]
把吹起 $\widetilde M$ 定义为并集
\[
  \widetilde M:=(M\setminus\{x_0\})\cup U_r\big/\sim,
\]
其中等价关系把点
$(z,[z_1:\cdots:z_n])\in U_r\setminus\mathbb{CP}^{n-1}$ 与
$\iota(z)\in M$ 识别。于是 $\widetilde M$ 是复流形，并且有嵌入
$\widetilde\iota:\mathbb{CP}^{n-1}\to\widetilde M$：它是零截面嵌入
\[
  \mathbb{CP}^{n-1}\longrightarrow U_r,
  \qquad [w_1:\cdots:w_n]\longmapsto(0,[w_1:\cdots:w_n])
\]
与包含 $U_r\lhook\joinrel\longrightarrow\widetilde M$ 的复合。其像是\emph{例外除子}
\[
  E:=\widetilde\iota(\mathbb{CP}^{n-1})\subset\widetilde M.
\]
此外，有全纯投影 $\pi:\widetilde M\to M$，满足 $E=\pi^{-1}(x_0)$，并且 $\pi$ 在
$\widetilde M\setminus E$ 上限制为到 $M\setminus\{x_0\}$ 的微分同胚。称
$(\widetilde M,\widetilde J)$ 为 $(M,J)$ 在 $x_0$ 处的\emph{复吹起}。可以证明，所得
流形的双全纯同构类与嵌入 $\iota$ 的选择无关。

需要强调的是，$M$ 上每个满足 $\iota^*J=J_0$ 的\BookRecordIndex{T0088}殆复结构 $J$，都会诱导
$\widetilde M$ 上唯一的\BookRecordIndex{T0088}殆复结构 $\widetilde J$，使投影 $\pi$ 全纯。复吹起的另一个有用
事实是，任意通过 $x_0$ 的 \BookRecordIndex{T0003}$J$-全纯曲线都能提升到吹起。因为这是局部陈述，只需考察
非常值 $J_0$-全纯映射 $u:(B^2,0)\to(\mathbb C^n,0)$ 的提升。每个这样的映射都形如
\begin{equation}\label{eq:9-3-1}
  u(z)=z^k\bigl(h_1(z),\ldots,h_n(z)\bigr),
  \qquad k\ge1,
\end{equation}
其中某个 $i$ 满足 $h_i(0)\ne0$；它提升为
\[
  \widetilde u(z)=
  \bigl(z^kh_1(z),\ldots,z^kh_n(z),[h_1(z):\cdots:h_n(z)]\bigr)\in L.
\]

若 $(M,J)$ 支持\BookRecordIndex{T0363}辛形式 $\omega$，则不难在吹起 $\widetilde M$ 上定义一个\BookRecordIndex{T0363}辛形式
$\widetilde\omega$，使它在例外除子 $E$ 中一条直线上的积分很小。不过，这种做法会掩盖
辛吹起与 $(M,\omega)$ 中嵌入球之间的紧密联系。McDuff \srcite{257} 指出，理解辛吹起
最具几何性的方式如下：切掉一个辛嵌入球
$B=\iota(B^{2n}(\rho))$ 的内部，再经由 Hopf 映射把边界球面
$\partial B\cong S^{2n-1}$ 压缩成例外除子 $E\cong\mathbb{CP}^{n-1}$。由此得到吹起的
另一种描述，即（适当光滑化的）商
\[
  \widetilde M_\rho:=M\setminus\iota(\operatorname{int}B)\big/\sim,
\]
其中，对 $z\in\partial B^{2n}(\rho)$ 和 $\lambda\in S^1$，令
$\iota(z)\sim\iota(\lambda z)$。在这个表述中，例外除子就是
$\iota(\partial B)$ 关于 $S^1$-作用的商。由于该 $S^1$-作用保持 $\omega$，形式
$\omega$ 可下降到此商；换言之，例外除子是相应的辛约化空间，也就是带有
Fubini--Study 形式 $\omega_{\mathrm{FS}}$ 的适当倍数的 $\mathbb{CP}^{n-1}$。因此，
吹起 $\widetilde M_\rho$ 从 $M$ 继承\BookRecordIndex{T0363}辛形式 $\widehat\omega_\rho$；它在
$M\setminus B$ 上与 $\omega$ 一致，在例外除子 $E$ 上限制为
$\rho^2\omega_{\mathrm{FS}}$，这里把 $\omega_{\mathrm{FS}}$ 归一化为在任意直线上的
积分等于 $\pi$。下面给出 $\widehat\omega_\rho$ 的精确公式。

辛吹起的这个描述揭示了它与嵌入球的联系，也就使它与\BookRecordIndex{T0120}非挤压定理的相关性变得可信。这个
思想已被形式化为\emph{\BookRecordIndex{T0360}辛切割}的概念，见 Lerman \srcite{237}。还要注意，这一过程可以
逆转：设 $(M,\omega)$ 中有一个辛嵌入的
$(\mathbb{CP}^{n-1},\rho^2\omega_{\mathrm{FS}})$ 副本 $E$，其法线丛可与\BookRecordIndex{T0414}重言线丛
识别；切掉 $E$ 的一个小邻域并粘入一个适当半径 $\rho+\varepsilon$ 的球，就得到一个
\BookRecordIndex{T0358}辛流形，称为 $(M,\omega)$ 沿 $E$ 的\emph{吹下}。（有些反直觉的是，吹起会减小
$M$ 的体积，而吹下则增大它。这个术语显然来自复范畴中的吹起操作；在复范畴中，吹起的确
把点 $x_0$ 扩大成了整个子流形。）

不过，上述定义不便于理解 $M$ 中的全纯曲线和其吹起中全纯曲线的关系，所以现在介绍
McDuff--Polterovich \srcite{263} 发展的另一种辛吹起方法。

\begin{proposition}\label{prop:9-3-3}
设 $(M,\omega)$ 是\BookRecordIndex{T0358}辛流形，$\iota:B^{2n}(r)\to M$ 是辛嵌入，而
$J\in\mathcal J(M,\omega)$ 是满足 $\iota^*J=J_0$ 的 $\omega$-\BookRecordIndex{T0348}相容\BookRecordIndex{T0088}殆复结构。则对每个
$\rho<r$，$(M,J)$ 在 $x_0:=\iota(0)$ 处的复吹起
$(\widetilde M,\widetilde J)$ 上都有\BookRecordIndex{T0363}辛形式 $\widetilde\omega_\rho$，满足：
\begin{enumerate}[label=(\roman*),leftmargin=3em]
\item 二形式 $\pi^*\omega$ 在
$\pi^{-1}\bigl(M\setminus\iota(B_r)\bigr)\subset\widetilde M$ 上与
$\widetilde\omega_\rho$ 一致；
\item $\widetilde\iota^*\widetilde\omega_\rho=\rho^2\omega_{\mathrm{FS}}$，即每条直线
上的积分为 $\rho^2\pi$ 的 Fubini--Study 形式；
\item $\widetilde\omega_\rho$ 与 $\widetilde J$ \BookRecordIndex{T0348}相容；
\item 若 $\Sigma$ 是闭定向二维流形，且
$\widetilde u:\Sigma\to\widetilde M$ 是光滑映射，则
\[
  \int_\Sigma(\pi\circ\widetilde u)^*\omega
  =\int_\Sigma\widetilde u^*\widetilde\omega_\rho
   +\pi\rho^2(\widetilde u\mathbin{\boldsymbol\cdot}E).
\]
\end{enumerate}
\end{proposition}

\begin{proof}
先证明\BookRecordIndex{T0414}重言线丛 $L\to\mathbb{CP}^{n-1}$ 零截面的一个邻域辛同胚于
$\mathbb R^{2n}$ 中的球壳。上面已经看到，$\delta$-球丛 $U_\delta$ 有坐标
$(z_1,\ldots,z_n,[z_1:\cdots:z_n])$。Fubini--Study 形式在齐次坐标中为
$\frac{i}{2}\partial\bar\partial\log|z|^2$
\sourcefootnote{原书把\BookRecordIndex{T0301}势函数印作 $\log|z^2|$；紧接着的公式及齐次坐标中的标准
Fubini--Study \BookRecordIndex{T0301}势函数均唯一确定此处应为 $\log|z|^2$。}，故二形式
\[
\begin{split}
  \omega_\rho
  &:=\frac{i}{2}\partial\bar\partial
       \bigl(|z|^2+\rho^2\log|z|^2\bigr)\\
  &=\frac{i}{2}\left(
      dz\wedge d\bar z
      +\rho^2\frac{dz\wedge d\bar z}{|z|^2}
      -\rho^2\frac{\bar z\mathbin{\boldsymbol\cdot}dz\wedge
                         z\mathbin{\boldsymbol\cdot}d\bar z}{|z|^4}
    \right)
\end{split}
\]
延拓为 $U_\delta$ 上的形式，并在例外除子上限制为
$\rho^2\omega_{\mathrm{FS}}$。这里记
\[
  dz\wedge d\bar z:=\sum_jdz_j\wedge d\bar z_j,
  \qquad
  \bar z\mathbin{\boldsymbol\cdot}dz:=\sum_j\bar z_j\,dz_j.
\]
现在考虑微分同胚
\[
  h_\rho:B^{2n}(\delta)\setminus\{0\}
  \longrightarrow B^{2n}\bigl(\sqrt{\rho^2+\delta^2}\bigr)
                   \setminus B^{2n}(\rho),
\]
其定义为
\[
  h_\rho(z):=\sqrt{|z|^2+\rho^2}\frac{z}{|z|}
  =\sqrt{1+\frac{\rho^2}{|z|^2}}\,z.
\]
一个略显繁琐但初等的计算表明，$h_\rho$ 对标准\BookRecordIndex{T0363}辛形式的拉回为
\[
  h_\rho^*\omega_0
  =\frac{i}{2}\left(
      dz\wedge d\bar z
      +\rho^2\frac{dz\wedge d\bar z}{|z|^2}
      -\rho^2\frac{\bar z\mathbin{\boldsymbol\cdot}dz\wedge
                         z\mathbin{\boldsymbol\cdot}d\bar z}{|z|^4}
    \right)
  =\omega_\rho.
\]
这个公式使我们能够精确描述吹起
$(\widetilde M_\rho,\widehat\omega_\rho)$。设
$\iota:B^{2n}(r)\to M$ 是辛嵌入。对任意 $\rho<r$，在
$M\setminus\iota(B^{2n}(\rho))$ 上令 $\widehat\omega_\rho=\omega$。这个形式可以光滑
延拓到例外除子上，因为可借助与 $h_\rho$ 相应的逆映射，把球壳
$\iota(B^{2n}(r)\setminus B^{2n}(\rho))$ 与
$U_\varepsilon\setminus\mathbb{CP}^{n-1}$ 识别。
\sourcefootnote{原书把识别映射印作 $h_\rho\circ\iota^{-1}$；按前一行给出的
$h_\rho$ 的源靶，这个复合映向欧氏球壳而非 $U_\varepsilon$，方向不合。此外，原书没有
定义这里的 $\varepsilon$。正文只保留证明所需的逆向识别，不擅自指定
$\varepsilon$ 的公式。}

接着构造微分同胚 $f:\widetilde M\to\widetilde M_\rho$，使拉回形式
$\widetilde\omega_\rho:=f^*\widehat\omega_\rho$ 满足命题要求。选取 $\delta>0$，使
\[
  \rho^2+\delta^2<(r-\delta)^2,
\]
并取光滑函数 $\beta:[0,r]\to[\rho,r]$，满足
\[
  \beta(s)=
  \begin{cases}
    \sqrt{s^2+\rho^2},&0\le s\le\delta,\\
    s,&r-\delta\le s\le r,
  \end{cases}
  \qquad
  0<s\le r\quad\Longrightarrow\quad0<\beta'(s)\le1.
\]
定义 $f:\widetilde M\to\widetilde M_\rho$ 为
\[
  f(\widetilde x):=
  \begin{cases}
    \pi(\widetilde x),
      &\pi(\widetilde x)\in M\setminus\iota(B^{2n}(r-\delta)),\\[2pt]
    \displaystyle\iota\left(\frac{\beta(|z|)z}{|z|}\right),
      &\pi(\widetilde x)=\iota(z),\quad0<|z|<r,\\[6pt]
    [\iota(w)],
      &\widetilde x=\widetilde\iota([w]),\quad[w]\in\mathbb{CP}^{n-1}.
  \end{cases}
\]
由构造，$\widehat\omega_\rho$ 在例外除子的一个删去邻域上的限制，可与球壳
$B^{2n}(\sqrt{\rho^2+\delta^2})\setminus B^{2n}(\rho)$ 上的标准形式
$\omega_0$ 识别。因此拉回形式
$\widetilde\omega_\rho:=f^*\widehat\omega_\rho$ 满足 (i) 和 (ii)。此外，直接计算可知，
只要对 $s>0$ 有 $\beta'(s)>0$，映射
$z\mapsto\beta(|z|)z/|z|$ 对 $\omega_0$ 的拉回便与 $J_0$ \BookRecordIndex{T0348}相容（见下面的
习题~\ref{ex:9-3-5}）。这说明 $\widetilde\omega_\rho$ 与 $\widetilde J$ \BookRecordIndex{T0348}相容。

为证明 (iv)，先设 $\widetilde u\mathbin{\boldsymbol\cdot}E=0$。由一般位置论证，还可
假设 $\widetilde u$ 横截于 $E$。沿 $E$ 中连接交数相反的交点的曲线作手术，可以在不改变
$\widetilde u$ 的同调类的前提下，使
$\widetilde u(\Sigma)\cap E=\varnothing$；不过，在消去交点的过程中，可能会减少
$\Sigma$ 的连通分支数，或增大某个连通分支的亏格。由此，$\widetilde u$ 与一个取值于
$M\setminus\iota(B^{2n}(r))\subset\widetilde M$ 的映射同调。对任何这样的映射，断言
(iv) 显然成立。这就证明了
$\widetilde u\mathbin{\boldsymbol\cdot}E=0$ 的情形。

现在令 $\widetilde u:\Sigma\to\widetilde M$ 是任意光滑映射，并记
$k:=\widetilde u\mathbin{\boldsymbol\cdot}E$。取光滑映射
$v:S^2\to\mathbb{CP}^{n-1}$ 表示同调类 $k[\mathbb{CP}^1]$，并记
$\widetilde v:=\widetilde\iota\circ v:S^2\to\widetilde M$。则
\[
  \widetilde v\mathbin{\boldsymbol\cdot}E=-k,
  \qquad
  \int_{S^2}\widetilde v^*\widetilde\omega_\rho=\rho^2\pi k,
  \qquad
  \int_{S^2}(\pi\circ\widetilde v)^*\omega=0.
\]
所以 $\widetilde u$ 与 $\widetilde v$ 的不交并和 $E$ 的交数为零。这就证明了 (iv)，也完成
命题~\ref{prop:9-3-3} 的证明。
\end{proof}

\begin{proof}[定理~\ref{thm:9-3-2} 的证明]
固定常数 $\varepsilon>0$，取面积形式 $\sigma\in\Omega^2(S^2)$，使
\[
  \int_{S^2}\sigma=\pi R^2+\varepsilon,
\]
并选取球 $B^2(R)$ 到 $S^2$ 的保面积嵌入。于是，题设给出的
$B^{2n}(r)$ 到 $B^2(R)\times V$ 的辛嵌入诱导辛嵌入
\[
  \iota:B^{2n}(r)\longrightarrow M:=S^2\times V.
\]
考虑由辛嵌入 $\iota$ 和常数 $\rho<r$ 按命题~\ref{prop:9-3-3} 确定的
$M$ 的吹起 $(\widetilde M,\widetilde\omega_\rho)$。令
\[
  A:=[S^2\times\{\mathrm{pt}\}]\in H_2(M;\mathbb Z),
\]
并用 $\widetilde A\in H_2(\widetilde M;\mathbb Z)$ 表示它到 $\widetilde M$ 中的提升。
因此，只要 $y\in V$ 满足
$x_0:=\iota(0)\notin S^2\times\{y\}$，类 $\widetilde A$ 就由子流形
$\pi^{-1}(S^2\times\{y\})$ 表示。

对每个 $J\in\mathcal J(M,\omega)$，$A$ 都是 \BookRecordIndex{T0002}$J$-不可分解的，故
引理~\ref{lem:7-1-8} 说明\BookRecordIndex{T0246}模空间 $\mathcal M(A;J)/G$ 对每个这样的 $J$ 都紧。选取
$M=S^2\times V$ 上的乘积\BookRecordIndex{T0088}殆复结构可知，对某个、因而对每个
$J\in\mathcal J_{\mathrm{reg}}$，\BookRecordIndex{T0283}求值映射
\[
  \operatorname{ev}:\mathcal M_{0,1}(A;J)\longrightarrow M
\]
的度数为一。因此
$\operatorname{GW}^{M}_{A,1}(\operatorname{PD}(\mathrm{pt}))=1$。由于\BookRecordIndex{T0246}模空间紧，
这意味着：对每个 $J\in\mathcal J(M,\omega)$ 和每个 $x\in M$，$M$ 中都有一条表示
类 $A$ 且通过 $x$ 的 $J$-全纯球面。

现在存在 $M$ 上一个 $\omega$-\BookRecordIndex{T0348}相容\BookRecordIndex{T0088}殆复结构 $J$，使 $\iota^*J$ 等于
$B^{2n}(r)$ 上的标准复结构 $J_0$。取 $J$-全纯球面
$u_0:S^2\to M$，使 $[u_0]=A$ 且 $x_0\in u_0(S^2)$。由
引理~\ref{lem:2-4-1}，集合 $Z_0:=u_0^{-1}(x_0)$ 有限。考虑映射
\[
  \widetilde u_0:=\pi^{-1}\circ u_0:
  S^2\setminus Z_0\longrightarrow\widetilde M.
\]
因为全纯曲线总能提升到复吹起，$\widetilde u_0$ 延拓为 $\widetilde M$ 中的
$\widetilde J$-全纯球面，仍记作 $\widetilde u_0$。由构造，延拓后的球面满足
$\widetilde u_0(Z_0)\subset E$。因此，由习题~\ref{ex:2-6-1}（\BookRecordIndex{T0186}交点正性的简单情形），
\[
  \widetilde u_0\mathbin{\boldsymbol\cdot}E>0.
\]
现在由命题~\ref{prop:9-3-3}(iv)，
\[
\begin{split}
  \pi R^2+\varepsilon
  &=\operatorname{Vol}(S^2)
   =\int_{S^2}u_0^*\omega
   =\int_{S^2}(\pi\circ\widetilde u_0)^*\omega\\
  &=\int_{S^2}\widetilde u_0^*\widetilde\omega_\rho
    +\pi\rho^2(\widetilde u_0\mathbin{\boldsymbol\cdot}E)
  \ge\pi\rho^2.
\end{split}
\]
因为这对每个 $\varepsilon>0$ 和每个 $\rho<r$ 都成立，所以 $R\ge r$。
\end{proof}

\begin{remark}\label{rem:9-3-4}
定理~\ref{thm:9-3-2} 的证明真正需要的是 $M=S^2\times V$ 上一个\BookRecordIndex{T0088}殆复结构 $J$，使
$\iota^*J$ 可积并被 $f_\beta^*\omega_0$ \BookRecordIndex{T0372}驯服，其中
\[
  f_\beta:B^{2n}(r)\setminus\{0\}
  \longrightarrow B^{2n}(r)\setminus B^{2n}(\rho),
  \qquad
  f_\beta(z):=\beta(|z|)|z|^{-1}z.
\]
一般而言，对任何这样的 $J$，投影 $S^2\times V\to S^2$ 都未必全纯。因此，我们施加了
$\pi_2(V)=0$，而不只是要求 $V$ \BookRecordIndex{T0058}半正。（若 $(M,\omega)$ \BookRecordIndex{T0058}半正，上述论证也成立。）若要把
证明推广到所有紧\BookRecordIndex{T0358}辛流形，就需要某种技术保证：对 $S^2\times V$ 上的每个\BookRecordIndex{T0088}殆复结构，
都有一条通过 $S^2\times V$ 中每个给定点、表示类 $[S^2\times\mathrm{pt}]$ 的稳定
\BookRecordIndex{T0003}$J$-全纯曲线。借助\BookRecordIndex{T0370}虚模循环可以做到这一点。
\end{remark}

\begin{exercise}\label{ex:9-3-5}
设 $f:(\mathbb C^n,0)\to(\mathbb C^n,0)$ 是形如
\[
  f(z)=\frac{\beta(|z|)}{|z|}z,
  \qquad \beta'(r)\ge0
\]
的微分同胚。证明 $\omega:=f^*\omega_0$ 与 $J_0$ \BookRecordIndex{T0348}相容。

\emph{提示：}对 $f$ 求微分，得到
\[
  df(z)\zeta=\beta'(|z|)\zeta_0+\frac{\beta(|z|)}{|z|}\zeta_1,
\]
其中 $\zeta_0$ 是 $\zeta$ 到 $\mathbb Rz$ 上的正交投影，且
$\zeta_0+\zeta_1=\zeta$。
\end{exercise}

到目前为止的讨论尚未涉及 $(M,\omega)$ 与其吹起
$(\widetilde M,\widetilde\omega)$ 的\BookRecordIndex{T0143}格罗莫夫--威滕不变量之间的关系。Hu \srcite{191}
利用 Li--Ruan \srcite{238} 和 Ionel--Parker \srcite{201} 的辛和公式，在非常一般的
情形解决了这个问题。用本书目前已有的工具很难表述十分一般的定理，因为\BookRecordIndex{T0058}半正流形 $M$
的吹起几乎从不\BookRecordIndex{T0058}半正，除非 $\dim M\le6$。下面的陈述留作习题。

\begin{exercise}\label{ex:9-3-6}
设 $(\widetilde M,\widetilde\omega)$ 是维数不超过六的\BookRecordIndex{T0358}辛流形 $(M,\omega)$ 的吹起。设
$a_1,\ldots,a_k\in H^*(M)$ 都是正次数类，并记
$\widetilde a_i:=\pi^*a_i\in H^*(\widetilde M)$。再给定
$A\in H_2(M)$，以循环 $Z_A\subset M\setminus\{x_0\}$ 表示；用
$\widetilde A\in H_2(\widetilde M)$ 表示 $\pi^{-1}(Z_A)$ 所代表的类。则
\[
  \operatorname{GW}^{\widetilde M}_{\widetilde A,k}
  (\widetilde a_1,\ldots,\widetilde a_k)
  =\operatorname{GW}^{M}_{A,k}(a_1,\ldots,a_k).
\]

\emph{提示：}考虑\BookRecordIndex{T0088}殆复结构 $J,\widetilde J$，其中 $J$ 在 $x_0$ 附近可积，且吹下映射
$\pi:(\widetilde M,\widetilde J)\to(M,J)$ 全纯。若 $C$ 是通过 $x_0$ 的
\BookRecordIndex{T0003}$J$-全纯曲线，则它到 $\widetilde M$ 的提升表示类 $\widetilde A-kL$，其中 $k$ 如
式~\eqref{eq:9-3-1}，而 $L$ 是例外除子 $E$ 中直线的类。因此，本习题归结为找到上述
\BookRecordIndex{T0088}殆复结构 $J$，使它对 $A$ 正则，并使对
$\operatorname{GW}^{M}_{A,k}(a_1,\ldots,a_k)$ 有贡献的 \BookRecordIndex{T0003}$J$-全纯曲线都不通过点 $x_0$。
\end{exercise}

\begin{exercise}\label{ex:9-3-7}
设 $Z$ 是 $(M,\omega)$ 中余维为 $2k$ 的辛子流形。明确构造 $M$ 沿 $Z$ 的法方向作辛
吹起的过程。

\emph{提示：}也许最简洁的做法，是把 $Z$ 的一个邻域与丛
$P\times_{U(k)}\mathbb C^k$ 中零截面的一个邻域识别，其中
$P\to Z$ 是主 $U(k)$-丛；然后以 $\mathbb C^k$ 在原点的吹起替换纤维
$\mathbb C^k$。见 Lerman \srcite{237}。
\end{exercise}

\paragraph{\BookRecordIndex{T0120}非挤压定理的推广。}

第一种推广，是考虑乘积流形
\[
  (M,\omega):=(S^2\times V,\pi_1^*\sigma+\pi_2^*\tau)
\]
中半径为 $r$ 的嵌入球，其中 $(V,\tau)$ 是任意\BookRecordIndex{T0358}辛流形。若 $V$ 闭，并假设任意闭\BookRecordIndex{T0358}辛流形
都有\BookRecordIndex{T0143}格罗莫夫--威滕不变量，则前面的论证无需本质修改便可成立，并给出预期的不等式
\eqref{eq:9-3-2}（见习题~\ref{ex:9-3-8}）。即使没有这一假设，Lalonde--McDuff
\srcite{223} 也通过一个相当复杂的几何构造证明了该不等式。这个构造适用于完全任意、
甚至非紧的\BookRecordIndex{T0358}辛流形 $V$。

第二种推广，是紧辛 $(2n-2)$-维流形 $V$ 上的辛二维球面丛
$S^2\lhook\joinrel\longrightarrow M\to V$。在这种情形，同样的技术适用；只要标准球
$(B^{2n}(r),\omega_0)$ 可辛嵌入 $(M,\omega)$，它们就证明不等式
\eqref{eq:9-3-2}，其中 $A$ 取纤维类。这个版本的\BookRecordIndex{T0120}非挤压定理在 Biran
\srcite{38,39} 的工作中起重要作用。

还可以考虑底为 $S^2$、纤维为 $(V,\tau)$ 的更一般哈密顿纤维丛。也就是说，假设存在
浸没 $\pi:M\to S^2$，使 $\omega$ 在每条纤维上的限制非退化，并辛同胚于
$(V,\tau)$。此时可以用比值
\[
  \operatorname{area}(M,\omega)
  :=\frac{\operatorname{Vol}(M,\omega)}{\operatorname{Vol}(V,\tau)}
  =\frac{\displaystyle\int_M\omega^n}
         {\displaystyle n\int_V\tau^{n-1}}
\]
度量 $(M,\omega,\pi)$ 的大小。若任意辛嵌入球的半径 $r$ 都满足
$\pi r^2\le\operatorname{area}(M,\omega)$，就称\BookRecordIndex{T0120}非挤压定理对
$(M,\omega,\pi)$ 成立。它不再在所有情形成立。不过，当纤维 $(V,\tau)$ \BookRecordIndex{T0355}辛非球面时，它
确实成立；证明及进一步结果见 McDuff \srcite{271}。这些思想的后续发展将在第~\ref{sec:9-6} 节
讨论。

也可以考虑更一般的底流形。不过，这个研究方向似乎希望不大，因为当 $\Sigma$ 是正亏格
黎曼曲面时，\BookRecordIndex{T0120}非挤压定理甚至并非对所有乘积
$(\Sigma\times V,\pi_1^*\sigma+\pi_2^*\tau)$ 都成立；见
\srcite{277}，习题 12.4。

\begin{exercise}\label{ex:9-3-8}
设 $(V,\tau)$ 是维数为 $2n-2$ 的\BookRecordIndex{T0358}辛流形，
$\sigma\in\Omega^2(S^2)$ 是\BookRecordIndex{T0363}辛形式，
$\iota:B^{2n}(r)\to S^2\times V$ 是辛嵌入，而
\[
  J\in\mathcal J(S^2\times V,\pi_1^*\sigma+\pi_2^*\tau)
\]
是满足 $\iota^*J=J_0$ 的\BookRecordIndex{T0088}殆复结构。假设存在一个表示同调类
$A:=[S^2\times\mathrm{pt}]\in H_2(S^2\times V;\mathbb Z)$、并通过点
$x_0:=\iota(0)$ 的 $J$-全纯\BookRecordIndex{T0340}稳定映射。证明
\begin{equation}\label{eq:9-3-2}
  \pi r^2\le\int_{S^2}\sigma.
\end{equation}
\end{exercise}

\begin{remark}[辛嵌入问题]\label{rem:9-3-9}
推广\BookRecordIndex{T0120}非挤压定理的另一个方向，是同时改变嵌入的定义域和靶。例如，可以考察嵌入多个等大
球、任意一组球、辛椭球或\BookRecordIndex{T0110}多圆盘。许多这类问题在四维已经解决，但人们对高维情形仍所知
甚少。对此感兴趣的读者，可以从论文 \srcite{65,167,175,278,231,56} 开始阅读。
\end{remark}
